% Generated draft: requires mathematical and source review.
\documentclass[11pt]{article}
\usepackage[margin=25mm]{geometry}
\usepackage{fontspec}
\setmainfont{Latin Modern Roman}
\usepackage{amsmath,amssymb,mathtools,unicode-math}
\setmathfont{Latin Modern Math}
\usepackage{enumitem,needspace,xurl,hyperref}
\setlength{\parindent}{0pt}
\setlength{\parskip}{5pt}
\setlength{\emergencystretch}{3em}
\newcommand{\N}{\mathbb{N}}
\newcommand{\Z}{\mathbb{Z}}
\newcommand{\Q}{\mathbb{Q}}
\newcommand{\R}{\mathbb{R}}
\newcommand{\C}{\mathbb{C}}
\newcommand{\F}{\mathbb{F}}
\newcommand{\A}{\mathbb{A}}
\newcommand{\PP}{\mathbb{P}}
\newcommand{\E}{\mathbb{E}}
\newcommand{\eps}{\varepsilon}
\newcommand{\dd}{\,\mathrm{d}}
\begin{document}
\section*{UnsolvedMath ID: 30000335}
\textbf{Source problem code:} catalogue entry
\subsection{Definitions and mathematical statement}
\textbf{Relative local-rigidity problem.} Prove that geometrically finite hyperbolic three-manifolds are locally rigid relative to the bending lamination on their convex-hull boundaries. In the standard complete orientable setting, write \(M=\mathbb H^3/\Gamma\), with \(\Gamma\le\mathrm{PSL}_2(\C)\) discrete and torsion free. Its convex core is the quotient of the hyperbolic convex hull of the limit set. When this core has the appropriate three-dimensional interior, its boundary is pleated and carries a measured bending lamination, whose transverse measure records bending angles. Fixing that measured lamination means fixing both its support and weights, not merely the underlying unmeasured lamination or induced boundary metric.\par A precise conditional deformation-space formulation fixes a marked topological core \(N\) and cusp data \(P\). Let \(\mathcal D\) be the intended space of marked geometrically finite structures, with the specified parabolic and topological conditions, modulo marking-compatible isometry, and let \(\beta\) record all boundary bending measured laminations using the marking. Local rigidity relative to bending means that for each permitted \(m_0\in\mathcal D\) there is a neighborhood \(U\) such that\[\{m\in U:\beta(m)=\beta(m_0)\}=\{m_0\}.\] This is isolation in the fixed-bending fiber, not necessarily a global uniqueness theorem, a claim that \(\beta\) is injective on all of \(\mathcal D\), or an infinitesimal rigidity statement. A representation-space version may use discrete faithful holonomy representations modulo conjugation and the corresponding algebraic topology, but the source does not state which deformation topology is intended.\par Essential scope conventions are missing. They include cusp behavior, preservation or variation of parabolic elements, the boundary marking, and treatment of elementary groups or convex cores with empty three-dimensional interior. In particular, Fuchsian cores are lower-dimensional and need a separate bending-boundary convention and possibly exclusion; they cannot simply be included under an undefined zero-bending map while asserting rigidity. No restriction to rational bending laminations, compact cores without cusps, or a fixed induced metric is inserted. The source's general wording is retained subject to recovering these hypotheses.\par The supplied triage checked 2026-08-21 labels the statement open and says it was not verified as solved in the stated generality. No precise known special case or theorem is supplied. This dated assessment requires independent verification and does not resolve the scope gaps above.
\subsection{Short English statement}
Nearby marked hyperbolic structures with exactly the same measured convex-core bending should coincide. The intended deformation space and cusp conventions are missing, and lower-dimensional Fuchsian cores require special treatment. Fixing only bending support would be a different condition.
\subsection{Sources}
\begin{itemize}\item Local UnsolvedMath registry, permanent ID 30000335, code catalogue entry; \emph{Relative Rigidity of Geometrically Finite Hyperbolic Three-Manifolds}, Oberwolfach Reports Open Problems.\item \emph{Low-Dimensional Manifolds (2006)}, Oberwolfach Reports, DOI 10.4171/owr/2005/45, pp. 2519--2570. \url{https://doi.org/10.4171/owr/2005/45}. Proposed year and year tag: 2005; no proposer named.\item Supplied literature label \emph{original Oberwolfach report}: \url{https://doi.org/10.4171/owr/2005/45}. No separate authoritative rigidity theorem or precise deformation-space citation is supplied.\item Literature triage checked 2026-08-21, status \texttt{open}; record updated 2026-09-16. Statement status \texttt{exact}, with confirmation reported without textual change; neither the confirmation nor the dated assessment was independently checked here.\end{itemize}
\end{document}
